% =====================================================================
\chapter{Genetic Algorithm}
% =====================================================================

\labinfo{13\_Genetic\_Algorithm}{genetic\_algorithm.py}{UCI Breast Cancer Wisconsin Diagnostic: 569
tumours, 30 numeric features}{Feature selection for a 5-NN classifier}{12 of 30 features selected;
CV accuracy 98.5\% (all 30: 96.8\%)}

\section{Dataset}
The \textbf{Breast Cancer Wisconsin Diagnostic} dataset (UCI) describes 569 tumours using 30 numeric
cell-nucleus measurements. The target is malignant or benign. I standardised the features and used a
stratified 70/30 split (398 training and 171 test tumours). The genetic algorithm selects a subset for
the 5-NN classifier while penalising unnecessarily large subsets.

\section{Code and experiments}
\begin{itemize}
\item \code{knn\_predict} is a compact 5-NN classifier (Lab~1) used inside the fitness function.
\item \code{Fitness} computes the penalised CV accuracy of a mask and stores every result in a dictionary
      keyed by the mask's bytes, so the same chromosome is never evaluated twice.
\item \code{tournament\_select}, \code{roulette\_select}, \code{one\_point\_crossover} and \code{mutate}
      are the genetic operators, each a few lines long.
\item \code{genetic\_algorithm} runs the algorithm's main loop and records the best and mean fitness of
      every generation.
\end{itemize}
\noindent Experiments: (1)~the main run with tournament selection and a second run with roulette-wheel
selection; (2)~five more runs with different random seeds, to see how stable the result is;
(3)~5-NN on the test set with all 30 features and with the selected ones.

\section{Results}

The full feature set has CV accuracy 0.9675. With tournament selection the best fitness reaches 0.9770
and the selected 12-feature subset has CV accuracy 0.9850. Roulette selection finds a 10-feature subset
with CV accuracy 0.9775. The penalty favours compact subsets while the cache avoids reevaluating identical
chromosomes.

\begin{table}[H]
\centering\small
\caption{Main GA runs (fitness and CV accuracy on the training set).}
\label{tab:ga-runs}
\begin{tabular}{@{}l c c c l@{}}
\toprule
\textbf{Run} & \textbf{Best fitness} & \textbf{CV accuracy} & \textbf{\#features} & \textbf{Selected features}\\
\midrule
All features (no GA) & 0.9475 & 0.9675 & 30 & all\\
GA, tournament & \textbf{0.9770} & \textbf{0.9850} & 12 & 1, 3, 8, 9, 10, 17, 20, 22, 24, 26, 27, 29\\
GA, roulette wheel & 0.9708 & 0.9775 & 10 & 8, 10, 12, 17, 20, 22, 23, 24, 27, 28\\
\bottomrule
\end{tabular}
\end{table}

\begin{table}[H]
\centering\small
\caption{Five more tournament runs with different random seeds.}
\label{tab:ga-seeds}
\begin{tabular}{@{}l c c c c c@{}}
\toprule
Seed & 1 & 2 & 3 & 4 & 5\\
\midrule
Features selected & 12 & 11 & 12 & 11 & 10\\
CV accuracy (train) & 0.985 & 0.980 & 0.983 & 0.980 & 0.978\\
Test accuracy & -- & -- & -- & -- & --\\
\bottomrule
\end{tabular}
\end{table}

Different seeds find different compact subsets with similar CV accuracy. This is typical for feature
selection: several measurements carry overlapping information, so multiple subsets can perform well.

\begin{table}[H]
\centering\small
\caption{5-NN on the test set (171 tumours), with all features and with the features chosen by the main GA run.}
\label{tab:ga-test}
\begin{tabular}{@{}l c S[table-format=1.4] S[table-format=1.4] S[table-format=1.4] S[table-format=1.4] l@{}}
\toprule
\textbf{Features} & \textbf{\#} & {\textbf{Accuracy}} & {\textbf{Macro P}} & {\textbf{Macro R}} & {\textbf{Macro F1}} & \textbf{Confusion (benign, malignant)}\\
\midrule
All features & 30 & 0.9708 & 0.9720 & 0.9650 & 0.9683 & --\\
GA-selected  & 12 & \bfseries 0.9766 & \bfseries 0.9770 & \bfseries 0.9750 & \bfseries 0.9760 & --\\
\bottomrule
\end{tabular}
\end{table}

On the training cross-validation score, the 12-feature subset raises the fitness from 0.9675 for all
features to 0.9770, while its measured subset CV accuracy is 0.9850. The selected subset keeps the
classifier compact without discarding the strongest tumour measurements.

\begin{figure}[H]
\centering
\includegraphics[width=0.48\linewidth]{\figpath{13\_Genetic\_Algorithm}{ga_convergence.pdf}}\hfill
\includegraphics[width=0.48\linewidth]{\figpath{13\_Genetic\_Algorithm}{selected_features.pdf}}
\caption{Genetic-algorithm convergence and selected feature masks.}
\end{figure}

\section{Source code}
\lstinputlisting[caption={\texttt{13\_Genetic\_Algorithm/genetic\_algorithm.py}}]{13_Genetic_Algorithm/genetic_algorithm.py}
\clearpage
